\subsection{A functional exchange with an inverse plateau}\label{sec:refined-mechanism}
The selected lower mechanism replaces the preceding construction. Its
parameters are selected jointly and verified exactly, so its revenue is integrated afresh
rather than obtained by adding gains from incompatible candidates.
The exchange threshold moves complete conditional menus. An upward jump
in that threshold induces a plateau in its generalized inverse, hence
changes a positive-measure interval of opposing reports.

Put
\[
 A=\frac23,\quad d=\frac12,\quad c=\frac{157}{500},\quad
 q=\frac{93}{500},\quad \nu=\frac{q^2}{2},\quad a=A+\nu,
\]
\[
 b=a+c=\frac{1496947}{1500000},\qquad
 s=a+d=\frac{1775947}{1500000},\quad
 \lambda=\frac45,\quad \bar u=\frac{3421}{10000}.
\]
The symbol $\bar u$ denotes a threshold parameter, not utility. Define
\begin{equation}\label{eq:refined-threshold}
 g(r)=\lambda(\bar u-r)\mathbf1_{[c,\bar u]}(r),\qquad
 h(r)=r+q+g(r),\qquad J=g(c)=\frac{281}{12500}.
\end{equation}
The increasing, right-continuous $h$ jumps upward at $c$. Its generalized
inverse on the needed reports is
\[
 k(t)=\begin{cases}
 t-q,&t\le d,\\
 c,&d<t\le d+J,\\
 (t-q-\lambda\bar u)/(1-\lambda),&d+J<t<\bar u+q,\\
 t-q,&t\ge\bar u+q.
 \end{cases}
\]
In particular, the proof uses $x<k(t)\Rightarrow h(x)<t$ and
$x\ge k(t)\Rightarrow h(x)\ge t$. The plateau does not satisfy
$h(k(t))=t$ at all its interior points.

\paragraph{Complete menus.}
For either bidder, let the opposing report be $w$, with
$t=\max(w)$, $\rho=\min(w)$ and $z=w_1+w_2$. In oriented menus,
``scarce'' means the opponent's high-valued physical item and ``safe''
means the other item. If the coordinates tie, designate item 2 as high.
Put $b_0=(4-\sqrt2)/3$, $T=1-2c/3$ and $T'=5/3-2c$.
Every menu includes empty allocation at zero payment.
\begin{enumerate}
\item On the closed region $Q=\{t\le A,z\le b\}$, if $t\le d$,
offer both singletons at $A$ and the bundle at $\max(b_0,z)$.
If $t>d$, set
\[
 C_0(k)=\frac56+\frac34k^2,\quad
 C=\max\{C_0(k(t)),z,k(t)+h(\rho)\},\quad
 B=\min\{A,C-k(t)\}.
\]
Offer safe at $B$, scarce at $A$, and bundle at $C$.
\item On the open region $E=\{A<t<T,\rho<c\}$, put
\[
 \delta=\frac9{16}(T-t)(T'-t),\quad \alpha=3t-2,
 \quad\beta=\frac{\alpha}{\alpha+2\delta}.
\]
Offer safe at $d+\delta$, scarce at $t$, bundle at $t+c+\delta$,
and allocation $(1,\beta)$ in scarce/safe order at price $t+\beta c$.
\item For every remaining opponent report, put
\[
 H=\max(0,w_1-a,w_2-a,z-b),\qquad
 P_j=H+\min(a,s-w_{3-j}),\qquad P_B=H+b.
\]
Offer both singletons and bundle at these prices, adding $g(\rho)$
to the safe singleton and bundle. The face $\rho=c$ receives the
positive fee in \eqref{eq:refined-threshold}.
\end{enumerate}
At zero maximal utility, retain empty only. At positive maximal utility,
retain every maximizing option. From the Cartesian product of the two
retained sets choose the first jointly feasible pair in lexicographic
order, bidder 1 first. Menu orders are empty/item 1/item 2/bundle for free
$Q$ and base, empty/safe/scarce/bundle for constrained $Q$, and
empty/safe/scarce/bundle/lottery for $E$.

\begin{proposition}[Complete functional-exchange mechanism]\label{prop:refined-feasible}
This selection is defined at every real report profile. It is Borel,
pointwise feasible, DSIC and ex-post IR. Its expected revenue is
\begin{equation}\label{eq:refined-revenue}
\begin{split}
 R_*={}&\frac{35791404252341621852527735747861637}
 {42525000000000000000000000000000000}
 +\frac{31}{1215}\sqrt2\\
 &-\frac{15059524650320123}{8305664062500000000000}\sqrt{493894}.
\end{split}
\end{equation}
\end{proposition}
\begin{proof}
The feasibility argument partitions opponent-menu pairs into base/base,
$Q$/base, $Q$/$Q$, $E$/base, $E$/$Q$ and $E$/$E$. The full all-real
inequalities are supplied in the source-bound structural proof identified
below; the following are the additional compatibility conditions created
by the changed parameters.
The base tariffs have strict bundle discount at least $q$. Raising the
safe and bundle prices by the same nonnegative fee therefore admits a
maximizing subset of an old base maximizer, including positive ties.
The base/base and base-facing comparisons can use that containment;
the complete new mechanism does not have predecessor containment.

In constrained $Q$, the bound $k+h(\rho)$ is redundant. On the fee
support its minimum sufficient slack is
\[
 [C_0(A-q)-(A-q)]-h(b-A)=\frac{37}{5000000}>0;
\]
outside the support the sum $z$ provides the required bound. Thus
$C=\max(C_0(k),z)$ everywhere, including reports with $C_0(k)>1$.
The cap at $A$ controls the bundle upgrade $\max(C-A,k)$.
The two generalized-inverse implications above establish the
complementary threshold comparisons at $Q$/base and $Q$/$Q$ junctions.

The $E$/$Q$ lottery comparison requires extra care because $a>A$.
A $Q$ type's lottery utility is at most $A-t+\beta(a-A)$.
On the open $E$ interval, $(t-A)/\beta=t-A+2\delta/3$ is strictly
increasing with infimum $q^2/2$. Since $a-A=q^2/2$, this utility
is strictly negative. The singleton and bundle comparisons and the
$E$/base and $E$/$E$ cases then leave a feasible maximizing pair,
as detailed in the finite case proof. The closed $Q$, open $E$ and
remaining base rule cover every boundary. In particular, the argument
does not assume that independently chosen tie priorities preserve supply.

The finite jointly feasible maximizer rule is Borel. For a fixed opponent,
every selected option maximizes the same menu; any deviation returns
another option of that menu. This proves DSIC for every report, and
empty proves IR. Payments are nonnegative and at most two by IR.
For each item, disjoint intervals of lengths $x_{1j},x_{2j}$ under a
uniform seed implement samplewise feasible supply. DSIC concerns expected
utility; truthfulness for each fixed seed is not asserted.

Revenue is computed by integrating each conditional menu over its own
unit square and then the full opponent square. Two implementations agree
on all three algebraic coefficients in \eqref{eq:refined-revenue}. One
uses rational polynomial dictionaries and simplex moments, the other
uses a direct $Q$ decomposition and Green boundary integrals over base
polygons. A third calculation in the present audit, importing no proof calculator,
independently reconstructs the full integral by vertical slicing and
recovers the same three coefficients. All include the negative base-to-$E$ correction on $A<t<a$
and the extra fee strip $c\le\rho\le b-A$, $A<t<b-\rho$. The inverse
plateau has positive opponent measure and is included. Ties do not alter
the integrals, while the selection argument above covers them pointwise.
\end{proof}

The new revenue exceeds the V4.6.1 predecessor by more than $17/2000000$.
The exact difference lies in
\[
 (0.00000855063547219779875,\ 0.00000855063547219779876).
\]
There is an ownership change throughout the open box of coordinate
radius $10^{-5}$ around
\[
 (v_1,v_2)=((103/200,1/10),(63/200,99/100)).
\]
The predecessor splits the goods; the new mechanism gives bidder 2 the
bundle. Uniform strict menu comparisons establish the entire box; its
sixteen corners alone would not prove a continuum claim.

\paragraph{Why these deformations help.}
In an auxiliary family with $\widetilde c=47/150$, $a=A=2/3$,
$d=1/2$ and $V=69/200$, integrating both bidders' full menu changes
produces a one-dimensional kernel
$\int_{\widetilde c}^V[\mathsf A(r)g(r)-\mathsf B(r)g(r)^2]\,dr$.
This calculation includes reports forced to move by the inverse plateau.
A second complete-menu variation explains the cap: in a proper
four-option menu with scarce price $A$, safe price $B$ and bundle price
$C$, $\partial_B\mathcal R=(2-3B)(C-B)$. If
$C-k=A+\eta>A$, capping the safe price at $A=2/3$ gains
$3k\eta^2/2+\eta^3/2>0$. These effects motivate globally admissible
changes. They do not establish unrestricted stationarity or optimality
of the selected parameters.

\subsection{Conditional optimality on the actual residual}\label{sec:wing-screening}
The selected mechanism has full randomized conditional screening
certificates for either bidder on three disjoint opponent regions,
$Q$, $E$ and
\[
 W_{\rm wing}=\{\max(w)\ge1-q,\ w_1+w_2\ge1+\bar u\}.
\]
Their areas are respectively
\[
 \frac{1746937679191}{4500000000000},\quad
 \frac{4867}{62500},\quad\frac{438867}{2500000},
\]
with union $2887322279191/4500000000000$, approximately $64.1627\%$
of the opponent square. This is conditional coverage, not a proportion
of the global revenue gap closed. The remaining area is unresolved,
not certified to have positive conditional regret.

The $Q$ theorem applies the preceding unrestricted envelope with the
capped price and verified actual top and anchor capacities. For $E$,
the earlier literal zero-capacity assertion on a junction face may fail
when $a>A$. Instead, zero scarce capacity in the adjacent open rectangle
$A<x<t$, $0<y<c$
forces every feasible competitor's utility to be horizontally constant.
Continuity and the two horizontal IC comparisons force its selected
scarce marginal to vanish on the required face. The contribution on $A<x<t$ of the singular line $y=c$ vanishes
by incentives; its $x>t$ term remains. This is a restriction imposed
by utility compatibility on that face.

\begin{theorem}[Screening on an occupied wing]\label{thm:wing}
Let $0<p\le1$, $1/3\le k\le1$, $z=k+p\ge4/3$, and
$(3k+1)p\ge k+1$. Suppose an actual pointwise residual admits the menu
$u^*(x,y)=\max(0,y-p,x+y-z)$, has first marginal zero on the open
strip $x<k$ and in the no-sale interior
$D_0=\{y<p,x+y<z\}$, and has second marginal zero on
$\{0\}\times[0,p)$. Then that menu is optimal over all randomized
DSIC/IR conditional mechanisms within the residual, with value
\[
 V(k,p)=pk(1-p)+(k+p)\left[(1-k)(1-p)+\frac{(1-k)^2}{2}\right].
\]
\end{theorem}
\begin{proof}
Let $u$ be the utility of an arbitrary complete measurable DSIC/IR
competitor, $a$ its selected allocation and $u_0=u(0,0)\ge0$.
The occupied left edge and horizontal strip force $u=u_0$ throughout
$D_0$ and $u(x,y)=u(0,y)$ on $x<k$. These are pointwise convex-utility
statements derived from chosen subgradients and Lipschitz continuity.
Averaging horizontal and vertical envelopes yields
\[
 R=\int\left[\frac{3x-1}{2}a_1+\frac{3y-1}{2}a_2\right]
 -\frac12\int u(0,y)\,dy-\frac12\int u(x,0)\,dx.
\]
The bottom trace is $u_0$, and
$\int u(0,y)dy=u_0+k^{-1}\int_{x<k}(1-y)a_2$.
Since $a=0$ almost everywhere on $D_0$, substitution gives
\begin{equation}\label{eq:wing-identity}
\begin{split}
 R+u_0={}&\int_G\left[\frac{3x-1}{2}a_1+\frac{3y-1}{2}a_2\right]\,dxdy\\
 &+\int_H\frac{(3k+1)y-(k+1)}{2k}a_2\,dxdy,
\end{split}
\end{equation}
where $G=\{x>k,x+y>z\}$ and $H=\{x<k,y>p\}$.
All displayed coefficients are nonnegative by the hypotheses.
The candidate saturates them, proving the value and equality.
\end{proof}

For $w=(t,\rho)$ in $W_{\rm wing}$, orient the scarce item first and
set $p=\min(\rho+q,1)$, $k=t+\rho-p$. The actual menu reduces to
safe at $p$ and bundle at $t+\rho$; an unaffordable scarce singleton
has no positive-utility demand. The opposite bidder occupies the scarce
item on the required strip and no-sale interior and both items on the
left-edge trace. These are strict comparisons for every maximizing option
in the required regions, across base, $Q$ and $E$ menus, including $p=1$.
The uniform sign margin is
\[
 (4-3q)(\bar u+q)-(2-q)=\frac{18601}{5000000}>0.
\]
Thus Theorem~\ref{thm:wing} applies to the actual residual for both bidders.
The source-bound proof includes the exceptional-face and orientation cases.
The theorem eliminates replacements with that residual held fixed.
Its rent elimination uses occupied sets and does not assert that the
remaining volume density alone supports arbitrary residual perturbations.

\paragraph{Stable source specification.}
The complete all-real case proof, independent integral derivation and
actual wing occupation proof are exported under
\begin{center}\small
\path{certificate/coordinated_primal_dual/source/V4_6_1_1_lower_bound/}
\end{center}
The proof notes in \path{research_log/} are
\path{refined_structure_audit.md}, \path{refined_revenue_audit.md} and
\path{residual_wing_screening.md}.

The exact implementation is
\path{verifier/refined_candidate.py}. The selected revenue is the single
expression \eqref{eq:refined-revenue}; earlier functional or reserve-only
experiments are not added to it.
